% Options for packages loaded elsewhere
\PassOptionsToPackage{unicode}{hyperref}
\PassOptionsToPackage{hyphens}{url}
\documentclass[
]{article}
\usepackage{xcolor}
\usepackage{amsmath,amssymb}
\setcounter{secnumdepth}{-\maxdimen} % remove section numbering
\usepackage{iftex}
\ifPDFTeX
  \usepackage[T1]{fontenc}
  \usepackage[utf8]{inputenc}
  \usepackage{textcomp} % provide euro and other symbols
\else % if luatex or xetex
  \usepackage{unicode-math} % this also loads fontspec
  \defaultfontfeatures{Scale=MatchLowercase}
  \defaultfontfeatures[\rmfamily]{Ligatures=TeX,Scale=1}
\fi
\usepackage{lmodern}
\ifPDFTeX\else
  % xetex/luatex font selection
\fi
% Use upquote if available, for straight quotes in verbatim environments
\IfFileExists{upquote.sty}{\usepackage{upquote}}{}
\IfFileExists{microtype.sty}{% use microtype if available
  \usepackage[]{microtype}
  \UseMicrotypeSet[protrusion]{basicmath} % disable protrusion for tt fonts
}{}
\makeatletter
\@ifundefined{KOMAClassName}{% if non-KOMA class
  \IfFileExists{parskip.sty}{%
    \usepackage{parskip}
  }{% else
    \setlength{\parindent}{0pt}
    \setlength{\parskip}{6pt plus 2pt minus 1pt}}
}{% if KOMA class
  \KOMAoptions{parskip=half}}
\makeatother
\setlength{\emergencystretch}{3em} % prevent overfull lines
\providecommand{\tightlist}{%
  \setlength{\itemsep}{0pt}\setlength{\parskip}{0pt}}
\usepackage{bookmark}
\IfFileExists{xurl.sty}{\usepackage{xurl}}{} % add URL line breaks if available
\urlstyle{same}
\hypersetup{
  hidelinks,
  pdfcreator={LaTeX via pandoc}}

\author{}
\date{}

\begin{document}

\section{Central Projection of a Vacuum Universe and Particle-like
States with a Spread Phase (Observational Note,
Draft)}\label{central-projection-of-a-vacuum-universe-and-particle-like-states-with-a-spread-phase-observational-note-draft}

\textbf{Author}: Noriaki Kihara \textbf{Affiliation}: WF System Co.,
Ltd.~/ Faculty of Engineering Science, Osaka University (graduate)
\textbf{ORCID}:
\href{https://orcid.org/0009-0004-6753-4020}{0009-0004-6753-4020}
\textbf{Version}: v1.0 (publication candidate, not peer-reviewed)
\textbf{Date}: June 2026 \textbf{License}: CC BY 4.0 \textbf{Concept
DOI}:
\href{https://doi.org/10.5281/zenodo.20543044}{10.5281/zenodo.20543044}
\textbf{Version DOI (v1.0)}:
\href{https://doi.org/10.5281/zenodo.20543045}{10.5281/zenodo.20543045}

\begin{center}\rule{0.5\linewidth}{0.5pt}\end{center}

\subsection{Character of this note}\label{character-of-this-note}

\textbf{This is an observational note in the form of a thought
experiment. It is not a proof paper or a claim paper.}

This note is a sequel to the preceding note {[}1{]} (the fifth frequency
and composite curvature). Whereas the preceding note observed that
imposing a sum-to-zero conservation condition on four equal-footing
frequencies makes a fifth direction \(\nu_5\) appear formally, and
developed the central projection with that as radius together with the
composite curvature radius, the present note observes---as an extension
of the same single thought experiment---that phase-bearing particle-like
states (phase-bearing local modes) can be placed on the subjective
geometry produced by that central projection, and that the spread
carried by such a particle can be characterized in the language of
frequency.

This note does not modify standard quantum theory, general relativity,
or special relativity. It proposes no new physical theory, proves no new
mathematical theorem, and changes no observable prediction. What it does
is to re-arrange and observe, from the single viewpoint of a
particle-like state, known geometric structures---dimensional reduction
by central projection {[}2,3{]}, standard facts of spherical and
Riemannian geometry {[}4,5{]}, and the fifth frequency / composite
radius of the preceding note {[}1{]}.

This note does not assert any of the following:

\begin{itemize}
\tightlist
\item
  any change to the mathematical predictions of the Standard Model,
  quantum theory, or relativity;
\item
  identification of a particle-like state with a concrete particle such
  as an electron, photon, or quark;
\item
  derivation of spin, statistics, gauge structure, or interactions;
\item
  any judgment of the stability of a local mode (equations of motion,
  energy functional, perturbative stability, etc.);
\item
  definition of an absolute position phase for a single particle;
\item
  endowing the post-projection 4-dimensional subjective geometry with a
  Lorentzian metric, causal structure, or physical spacetime;
\item
  prediction of values of new physical constants or the cosmological
  constant, or any new observable prediction;
\item
  proof of any new mathematical theorem.
\end{itemize}

Evaluation and interpretation are left to the reader.

\begin{center}\rule{0.5\linewidth}{0.5pt}\end{center}

\subsection{Abstract}\label{abstract}

Taking an exterior-free vacuum universe as a single intrinsic geometric
system, we observe that particle-like states with a spread phase
(phase-bearing local modes) can be placed on the subjective geometry
produced by its central projection. Starting from the fifth frequency
\(\nu_5\) and composite radius \(R_{\mathrm U}=1/\nu_5\) introduced in
the preceding note {[}1{]}, and from the central projection {[}2,3{]},
we observe the following. (i) We define the post-projection subjective
geometry as the homogeneous 4-sphere \(S^4(R_{\mathrm U})\) in
\(\mathbb{R}^5\) (sectional curvature \(1/R_{\mathrm U}^2\)), and regard
the central projection as a gnomonic chart covering one hemisphere. We
do not endow it with a Lorentzian metric or causal structure, and do not
call it physical spacetime. (ii) Since this geometry is homogeneous,
with no origin and no reference axis, no absolute position phase is
defined for a single particle. (iii) A particle-like state is
represented minimally by \(P_a=(\boldsymbol{\nu}_a,\,R_a)\) as a
phase-bearing local mode on this geometry. (iv) Under the model
convention that reads a frequency as the winding number of phase and
reads the spread as an effective diameter, the spread (full width) of a
particle is read as twice the curvature radius,
\(W_a=2R_a=2/\nu_{R,a}\), and the spread frequency coincides with the
curvature frequency (\(\nu_w=\nu_{R,a}\)). The scope is limited to
defining the single-particle geometric stage; the relative phase,
interference, and statistics of multiple particles are deferred to a
subsequent note.

\begin{center}\rule{0.5\linewidth}{0.5pt}\end{center}

\subsection{§1 Basic assumptions}\label{basic-assumptions}

First, consider a vacuum universe with no exterior space.

Ordinary intuition tends to imagine the universe as a small region
inside some exterior space. Here, however, we assume no such exterior
space. There is no space outside the universe. Hence the early universe
has no position or size measured from outside.

Then the vacuum universe is not an object with external coordinates, but
is defined only as an intrinsic geometric structure.

We represent its minimal model by an intrinsic frequency vector

\[\boldsymbol{\nu}=(\nu_1,\nu_2,\nu_3,\nu_4)\]

and the composite radius determined from it,

\[\nu_5^2=\nu_1^2+\nu_2^2+\nu_3^2+\nu_4^2\]

(preceding note {[}1{]}).

Here \(\nu_i\) is not an ordinary wavelength or a coordinate on an
external space, but a discrete, formal frequency parameter describing
the intrinsic geometric system. In this note we treat \(\nu_i\) as a
formal (dimensionless) parameter carrying no external dimension of time
or length (\([T^{-1}]\) or \([L]\)). Consequently \(x_i=1/\nu_i\) and
the composite radius \(R_{\mathrm U}\) determined from it also carry no
real length in external units; they have meaning only as intrinsic scale
ratios of the system.

Corresponding to the composite radius, we define the system-wide
central-projection radius

\[R_{\mathrm U}=\frac{1}{\nu_5}.\]

The subscript U indicates that this is a quantity of the system (the
vacuum universe) as a whole, to be distinguished from the per-particle
local radius \(R_a\) (§4) introduced later.

\begin{center}\rule{0.5\linewidth}{0.5pt}\end{center}

\subsection{§2 Central projection}\label{central-projection}

On this intrinsic geometric system we define a central projection
{[}2,3{]}.

For geometric computation, set the coordinates of each axis to

\[x_i=\frac{1}{\nu_i}\qquad(i=1,2,3,4,5).\]

This does not mean a physical wavelength or a real coordinate on an
external space; it is merely an auxiliary coordinate for formulating the
central projection.

The fifth component is generated by the same rule from \(\nu_5\), which
determines the composite radius:

\[x_5=\frac{1}{\nu_5}=R_{\mathrm U}.\]

Thus the point

\[x=\left(\frac{1}{\nu_1},\ \frac{1}{\nu_2},\ \frac{1}{\nu_3},\ \frac{1}{\nu_4},\ \frac{1}{\nu_5}\right)\]

has all five components determined by the single rule \(x_i=1/\nu_i\).
For this point we define the central projection onto the 4-sphere
\(S^4(R_{\mathrm U})\) of radius \(R_{\mathrm U}\) in \(\mathbb{R}^5\)
as

\[x'=\frac{R_{\mathrm U}}{|x|}\,x,\qquad |x|=\sqrt{\textstyle\sum_{i=1}^{5}x_i^2}.\]

(see {[}2,3{]}). Since \(x_5=R_{\mathrm U}\) is fixed system-wide, the
source point always lies on the hyperplane \(x_5=R_{\mathrm U}\). This
is the hyperplane tangent to \(S^4(R_{\mathrm U})\) at the north pole
\(N=(0,0,0,0,R_{\mathrm U})\), and the central projection from the
origin (the sphere's center) is precisely the gnomonic projection (here
used as a Beltrami-type projective chart) {[}2,3{]}. Its image is
exactly the open hemisphere on the side of the point of tangency
(\(x'_5>0\)); moving over all admissible parameters, the image never
exceeds this single hemisphere.

Here \(\nu_i\) is treated not as an ordinary positive frequency but as a
formal parameter carrying a sign or phase direction. Consequently the
auxiliary coordinate \(x_i=1/\nu_i\) may take both positive and negative
values on the local chart, covering the whole tangent hyperplane (the
whole hemisphere). If \(\nu_i\) were restricted to positive values, the
image of the gnomonic chart would be restricted not to the whole
hemisphere but to a sub-region of it (the positive-octant side).

Accordingly, in this note we define the subjective geometry not as ``the
image of the map'' (a hemisphere) but as the manifold
\(S^4(R_{\mathrm U})\) itself. The above gnomonic projection is
positioned as one coordinate chart covering this manifold (a projective
patch over one hemisphere). In gnomonic projection, straight lines on
the tangent hyperplane map to great circles (geodesics) on the
sphere---this is the defining property of a projective chart, and it
connects directly to the discussion of \(\sigma_R\) in the preceding
group of notes {[}2,3{]}. Note that gnomonic projection is not an
isometry. Moreover, since a single gnomonic chart covers only one
hemisphere of \(S^4(R_{\mathrm U})\) rather than the whole sphere, an
atlas of several such charts is needed to treat the whole sphere. In
this note we take the stance that global statements (homogeneity,
absence of absolute position, etc.) are made on the manifold
\(S^4(R_{\mathrm U})\), while explicit coordinate computations are
carried out on a single chart.

The auxiliary coordinate \(x_i=1/\nu_i\) introduced here does not invert
the metric of \(\boldsymbol{\nu}\) space as it stands. Hence the
Euclidean norm in \(\boldsymbol{\nu}\) space and the Euclidean norm in
\(x\) space are not identified. In particular, the fifth coordinate
\(x_5=1/\nu_5=R_{\mathrm U}\) does not in general equal
\(\sqrt{x_1^2+x_2^2+x_3^2+x_4^2}\). In this note we use \(x_i\) only as
an auxiliary coordinate for describing the central projection.

By this construction, a subjective manifold \(S^4(R_{\mathrm U})\) is
determined from the intrinsic geometric system. The subjective geometry
here is not a background space given from outside, but an intrinsic
manifold defined by the construction of the central projection. On it we
adopt, as a model choice, the round (standard) metric induced by the
auxiliary embedding \(\mathbb{R}^5\supset S^4(R_{\mathrm U})\).

\begin{center}\rule{0.5\linewidth}{0.5pt}\end{center}

\subsection{§3 The post-projection subjective
geometry}\label{the-post-projection-subjective-geometry}

The subjective geometry \(S^4(R_{\mathrm U})\) is the 4-sphere of radius
\(R_{\mathrm U}\) in \(\mathbb{R}^5\), and its intrinsic dimension is 4.
In this note we adopt, as a model, the round (standard) metric induced
by the auxiliary embedding (§2).

In this subjective geometry, for a sufficiently small geodesic triangle,
the sum of the interior angles exceeds 180 degrees (the angular excess
is positive). This shows that this intrinsic geometry has positive
sectional curvature. Read as the standard spherical geometry {[}4{]} of
radius \(R_{\mathrm U}\), the sectional curvature is
\(1/R_{\mathrm U}^2\).

We do not introduce a Lorentzian metric or causal structure on this
4-dimensional structure {[}5{]}. Hence what this note treats is not
physical spacetime itself, but a positively curved 4-dimensional
subjective geometry at a stage prior to a spacetime interpretation. The
emergence of the metric signature and the time direction is left to the
sum-to-zero framework of the preceding note {[}1{]}, and is not entered
into here. \(S^4(R_{\mathrm U})\) may be regarded as the
constant-curvature space corresponding to the Euclideanization of the
Lorentzian de Sitter space \(dS_4\) (sectional curvature
\(+1/R_{\mathrm U}^2\)), consistent with this note's policy of
introducing no causal structure.

On the other hand, the direction of the axis that produces this
curvature---i.e., the direction appearing outside the sphere in the
auxiliary embedding representation---cannot be known from inside the
system (this note does not assume an exterior space as real; this is
merely a direction in the auxiliary embedding representation). The
curvature is symmetric in all four directions, and there is no special
direction.

As a manifold, \(S^4(R_{\mathrm U})\) is a homogeneous space (the
rotation group \(O(5)\) acts transitively on points), and has no
distinguished origin or reference axis. The north pole \(N\) and the
equator that appear in the gnomonic chart of §2 are artifacts specific
to that chart, not properties of the manifold itself. Therefore no
absolute position is defined on \(S^4(R_{\mathrm U})\). Position is not
a quantity given by external coordinates; it has meaning only as a
relation between local modes (a coordinate-invariant quantity such as
geodesic distance) (§6). In this note we do not define a position phase
for a single local mode.

A particle-like state is then not a point placed in an external space,
but a local mode with a spread phase appearing on the centrally
projected subjective geometry (the meaning of ``spread phase'' is fixed
in §4 and §5).

\begin{center}\rule{0.5\linewidth}{0.5pt}\end{center}

\subsection{§4 Definition of a particle-like
state}\label{definition-of-a-particle-like-state}

In this note we do not assume a particle in advance as a physical
entity.

A particle-like state is defined as the following geometric object:

\[\text{particle-like state}\ :=\ \text{a local mode with a spread phase appearing on the post-projection subjective geometry}.\]

Here a ``local mode'' does not mean an eigenmode whose stability has
been determined. The judgment of stability (equations of motion, energy
functional, perturbative stability, etc.) is not made in this note and
is deferred to a subsequent note. This note merely defines the geometric
stage on which such local modes are placed.

Moreover, ``phase-bearing'' here does not mean an absolute position
phase on the subjective geometry. The phase that this note treats for a
single particle-like state is the intrinsic periodic structure
corresponding to the spread \(W_a\), i.e.~the spread phase (§5). The
absolute position phase \(\phi_a\) is not defined for a single particle
(§6).

Here a local mode means a local configuration supported (with support)
on \(S^4(R_{\mathrm U})\); its origin is left open, and it is
characterized by the intrinsic frequency \(\boldsymbol{\nu}_a\) and the
curvature radius \(R_a\). In this note we do not fix this configuration
as any specific mathematical entity such as a local deformation of the
metric or a wave packet of a field. The identification of what it is on
the subjective geometry (a wave packet of a scalar field, a local
deformation of the curvature, etc.) is deferred, together with the
introduction of dynamics, to a subsequent note.

In this note we treat only a single particle-like state. In this case we
do not define an absolute position phase for the particle-like state. In
this note's setting, which assumes no exterior space, position is not a
quantity given by external coordinates; it has meaning only as a
relation between several particle-like states.

Therefore the particle-like state treated in this note is represented
minimally by

\[P_a=(\boldsymbol{\nu}_a,\ R_a),\]

where \(\boldsymbol{\nu}_a\) is the intrinsic frequency structure and
\(R_a\) is the composite radius (local projection radius) of the
particle-like state. \(R_a\) is a per-particle quantity, to be
distinguished from the system-wide projection radius \(R_{\mathrm U}\)
of §1 (the two do not in general coincide).

The position phase \(\phi_a\) and the relative phase \(\Delta\phi_{ab}\)
become definable only when at least two particle-like states \(P_a,P_b\)
exist. This note does not enter into that formulation (§6).

At this stage, however, we do not identify this with a concrete particle
such as an electron, photon, or quark. What this note observes is
limited to the geometric structure that a phase-bearing local state can
be placed by the central projection of a vacuum universe.

\begin{center}\rule{0.5\linewidth}{0.5pt}\end{center}

\subsection{§5 Spread phase and the size of a particle (model
convention)}\label{spread-phase-and-the-size-of-a-particle-model-convention}

Up to the previous section, a particle-like state was defined as a
phase-bearing local mode on the subjective geometry. In this section we
characterize the ``spread'' carried by this local mode in the language
of frequency. This is an observation applying the viewpoint of frequency
\(\nu\)---which the preceding note {[}1{]} adopted as its basic quantity
of description---to the size of a particle. \textbf{What follows is not
a derivation but a model convention adopted in this note.}

First we fix one convention for the correspondence between phase and
frequency. In this note we read the frequency \(\nu\) as the number of
times the phase makes one full turn (360 degrees, \(2\pi\)); that is,

\[\nu=1\ \longleftrightarrow\ \text{phase }360^\circ\,(=2\pi).\]

This convention serves to associate phase and frequency directly,
without going through wavelength. If wavelength is used as the
intermediary, the coefficient wavers depending on whether \(1/\lambda\)
or \(2\pi/\lambda\) is called the frequency; hence this note takes the
winding number of phase as basic.

Furthermore, the ``spread'' of a particle in this note denotes not the
arc length on the phase circle but the \textbf{effective diameter}
occupied by the phase mode. That is, we read the geometric quantity
corresponding to one full turn of phase not as the circumference
\(2\pi R_a\) but as the effective width from center to both ends,
\(2R_a\). This is a model convention of this note, not an identification
with the ordinary circumference \(2\pi R_a\). There are two reasons for
adopting this reading. First, to express a particle-like state as a
centrally symmetric occupied range, the full width from center to both
ends is the natural measure. Second, taking the arc length
(circumference) would introduce a factor \(2\pi\), and we wish to avoid
a coefficient stacked on top of this note's framework, which treats the
phase as a winding number (\(\nu=1\leftrightarrow360^\circ\)). Under
this convention we associate the spread (full width) of a particle with
one full turn of phase, and its half (from center to an end) with a half
turn, and we write the frequency accompanying this spread as \(\nu_w\).

On the other hand, writing the frequency in the composite-radius
(curvature) direction of the particle-like state as \(\nu_{R,a}\), by
the same rule as in §1 and §4 we have

\[R_a=\frac{1}{\nu_{R,a}}.\]

In this note's convention, the curvature radius \(R_a\) is associated
with the effective radius (a half turn) of the phase circle.

Here \(\nu_w\) was originally introduced as an independent quantity, the
frequency accompanying the spread. In this note we place a model
identification regarding the pulsation of the spread and the pulsation
in the curvature direction as \textbf{one and the same intrinsic
vibration}. That is, not as an equality by definition but as an
assumption that the two vibrations are the same,

\[\boxed{\ \nu_w=\nu_{R,a}\ }\]

is set. Hereafter we do not retain \(\nu_w\) as an independent quantity
but unify on \(\nu_{R,a}\). Then, from the relation between the
curvature radius \(R_a\) (effective radius) and the full width of the
spread (effective diameter),

\[\boxed{\ W_a\equiv(\text{spread of the particle; full width})=2R_a=\frac{2}{\nu_{R,a}}\ }\]

can be read off. That is, in this note we take the curvature radius
\(R_a\) to be half the particle's spread (the effective radius), and the
full width of the spread to be its effective diameter \(2R_a\). This is
merely a re-reading between radius and diameter at the same frequency
\(\nu_{R,a}\), and no independent degree of freedom other than the
factor 2 enters between the two.

This convention assigns an effective spread to the particle-like state.
A particle is neither a point nor a region on a background space, but a
phase mode of spread \(W_a=2R_a=2/\nu_{R,a}\) whose effective radius is
the curvature radius \(R_a\). The larger the curvature frequency
\(\nu_{R,a}\), the smaller the spread; the smaller \(\nu_{R,a}\), the
larger the spread. If \(R_a\) is interpreted as a length on the geodesic
scale of the subjective geometry \(S^4(R_{\mathrm U})\), its upper bound
is naturally given by the geodesic extent of \(S^4(R_{\mathrm U})\) (for
example, the antipodal distance \(\pi R_{\mathrm U}\)). If \(R_a\) is
assumed to have this global scale as an upper bound, then in the limit
where \(\nu_{R,a}\) takes its minimal value the spread of the particle
approaches the scale of the whole subjective geometry. However, in this
note we leave it undetermined whether this upper bound is imposed as a
model convention, and confine ourselves to noting the possibility.

This section stops here. When several particles with a spread exist, how
their relative positions are converted into frequency, and the
interference, the discretization of placement, and the statistical
structure that then appear, require treating the relative phase and go
beyond the single-particle framework. These are made the subject of a
subsequent note; in this note we confine ourselves to observing, up to
the single-particle model convention that a particle has the spread
\(W_a=2R_a\).

\begin{center}\rule{0.5\linewidth}{0.5pt}\end{center}

\subsection{§6 One-particle universe and many-particle
universe}\label{one-particle-universe-and-many-particle-universe}

When there is only one particle-like state, the only things that can be
defined are the intrinsic frequency structure \(\boldsymbol{\nu}_a\) and
the spread \(W_a=2R_a\). Here \(W_a\) is not a ``region in which it is
placed'' but an intrinsic scale (effective diameter) of the
particle-like state. Position is not defined, because there is no
counterpart to measure it against. In a one-particle universe, no
absolute position phase exists. That it has a spread but no position is,
on a homogeneous space, rather a natural consequence.

Only when several particle-like states

\[P_a,\quad P_b\]

exist does the relative phase between them

\[\Delta\phi_{ab}\]

come into question. This relative phase is a candidate quantity for
formulating distance, direction, and relations in the subjective
geometry. Since the system has no origin and no reference axis, absolute
position is not defined, and only the relations between particles carry
physical meaning.

However, on the curved subjective geometry, a simple subtraction of
phases is not in general coordinate-invariant. To define the relative
phase / relative position globally and coordinate-invariantly requires
additional formulation such as the choice of local chart, gauge fixing,
and parallel transport. This note does not carry out that formulation.
The relative phase, relative position, interference, discretization of
placement, and statistical structure in a many-particle system are made
the subject of a subsequent note.

Therefore in this note we confine ourselves to reading space not as an
external container in which particles are placed, but as the relations
among several phase-bearing local states.

\begin{center}\rule{0.5\linewidth}{0.5pt}\end{center}

\subsection{§7 Relation between the vacuum universe and the
particle-like
state}\label{relation-between-the-vacuum-universe-and-the-particle-like-state}

The vacuum universe is a single intrinsic geometric system with no
exterior. Its central projection makes the subjective geometry appear.

A particle-like state is a phase-bearing local mode appearing on that
subjective geometry.

Hence a particle-like state is not placed directly in the original
constituting geometric space, nor is it brought into the vacuum universe
from an exterior space.

Rather, a particle-like state appears within the map

\[\text{central projection of the vacuum universe}\ \rightarrow\ \text{subjective geometry}\ \rightarrow\ \text{phase-bearing local mode}.\]

In this sense, a particle-like state is a local mode appearing on the
map structure that holds between the pre-projection geometry and the
post-projection subjective geometry.

\begin{center}\rule{0.5\linewidth}{0.5pt}\end{center}

\subsection{§8 What this note does not
treat}\label{what-this-note-does-not-treat}

This note does not treat the following:

\begin{itemize}
\tightlist
\item
  judgment of the stability of a local mode (equations of motion, energy
  functional, perturbative stability);
\item
  definition of an absolute position phase for a single particle;
\item
  endowing the post-projection geometry with a Lorentzian metric, causal
  structure, or physical spacetime;
\item
  the origin of spin, bosonic / fermionic character;
\item
  gauge particles, electromagnetic interaction, gravitational
  interaction;
\item
  composite particles, quarks, generation structure, the Higgs field;
\item
  modification of the Standard Model, new observable predictions;
\item
  conversion of the relative position between several particles into
  frequency, interference, and statistical structure.
\end{itemize}

These are extension topics that could be examined if the geometric
structure of a post-projection phase-bearing particle-like state were
further developed. At the present stage, however, they go beyond the
scope of this note.

\begin{center}\rule{0.5\linewidth}{0.5pt}\end{center}

\subsection{§9 Scope of the claims of this
note}\label{scope-of-the-claims-of-this-note}

The claims of this note are limited to the following range.

First, an exterior-free vacuum universe can be posited as a single
intrinsic geometric system.

Second, a central projection can be defined on that intrinsic geometric
system.

Third, the central projection determines the subjective geometry as the
homogeneous 4-sphere \(S^4(R_{\mathrm U})\) in \(\mathbb{R}^5\)
(sectional curvature \(1/R_{\mathrm U}^2\)), and gives a gnomonic chart
covering one hemisphere of it. An atlas of several charts is required to
treat the whole sphere. No Lorentzian metric or causal structure is
introduced here.

Fourth, a particle-like local mode can be placed on that subjective
geometry.

Fifth, this particle-like state is observed not as a point in an
external space but as a phase-bearing local mode on the centrally
projected subjective geometry, minimally as
\(P_a=(\boldsymbol{\nu}_a,R_a)\).

Sixth, under the model convention that reads a frequency as the winding
number of phase and the spread as an effective diameter, the spread
(full width) of a particle-like state can be read as twice the curvature
radius, \(W_a=2R_a=2/\nu_{R,a}\), and the spread frequency coincides
with the curvature frequency (\(\nu_w=\nu_{R,a}\)).

Seventh, an absolute position phase is not defined for a single
particle, and a relative phase comes into question only for two or more
particle-like states. Its formulation is not carried out in this note.

\begin{center}\rule{0.5\linewidth}{0.5pt}\end{center}

\subsection{Summary}\label{summary}

The central observation of this note can be summarized in one sentence.

\begin{quote}
We observe a particle not as a point-like object placed in an external
space, but as a local mode with a spread phase appearing on the
subjective geometry produced by the central projection of an
exterior-free vacuum geometric system.
\end{quote}

At this stage, this note does not attempt to complete a theory of
particles. What it does is to define the geometric stage on which a
particle-like state appears, and to set the convention characterizing
the spread carried by such a particle in terms of frequency.

That is, it presents, as an observational note, the map structure

\[\text{vacuum universe}\ \rightarrow\ \text{central projection}\ \rightarrow\ \text{subjective geometry}\ \rightarrow\ \text{phase-bearing particle-like state (spread }W_a=2R_a\text{)}.\]

\begin{center}\rule{0.5\linewidth}{0.5pt}\end{center}

\subsection{References}\label{references}

{[}1{]} N. Kihara (2026). \emph{A thought experiment on a model particle
(continued)---the fifth frequency and composite curvature seen from the
sum-to-zero condition}. Observational note (same author; preceding
note). Concept DOI:
\href{https://doi.org/10.5281/zenodo.20535894}{10.5281/zenodo.20535894}.

{[}2{]} N. Kihara (2026). \emph{Definition of the spherical projection
\(\sigma_R\) (radial projection)---the foundational map of the
central-projection series}. Observational note (same author). Concept
DOI:
\href{https://doi.org/10.5281/zenodo.20462569}{10.5281/zenodo.20462569}.

{[}3{]} N. Kihara (2026). \emph{A geometric formulation of 4-dimensional
space by central projection}. Observational note (same author). Concept
DOI:
\href{https://doi.org/10.5281/zenodo.19427780}{10.5281/zenodo.19427780}.

{[}4{]} M. P. do Carmo (1992). \emph{Riemannian Geometry}. Birkhäuser.
(spherical geometry, sectional curvature, angular excess of geodesic
triangles)

{[}5{]} B. O'Neill (1983). \emph{Semi-Riemannian Geometry, with
Applications to Relativity}. Academic Press. (distinction between
Riemannian and Lorentzian geometry, causal structure)

(Note: external references use the Concept DOI.)

\begin{center}\rule{0.5\linewidth}{0.5pt}\end{center}

\subsection{Revision history}\label{revision-history}

\begin{itemize}
\tightlist
\item
  \textbf{v1.0 (2026-06)}: Publication candidate. Reflecting a sixth
  round of external review, the subjective geometry is defined as the
  manifold \(S^4(R_{\mathrm U})\) itself (not as the hemispherical image
  of the map), with the central projection positioned as a gnomonic
  chart covering one hemisphere; homogeneity (\(O(5)\) transitivity) and
  the absence of absolute position thereby follow as properties of the
  manifold (the north pole and equator are chart artifacts). Also:
  commitment to the round metric stated as a model choice (gnomonic is
  not an isometry); correspondence to the Euclideanization of \(dS_4\)
  noted; \(\nu_i\) treated as a dimensionless formal parameter carrying
  a sign/phase direction; \(\nu_w=\nu_{R,a}\) clarified as a model
  identification rather than a definition, unifying on \(\nu_{R,a}\);
  \(R_a\) interpreted as a geodesic-scale length on
  \(S^4(R_{\mathrm U})\); carrier of the local mode (a configuration
  supported on \(S^4\)) stated; need for an atlas to cover the whole
  sphere stated; standard geometry references {[}4,5{]} added.
  (Translated from the Japanese v1.0; see the Japanese version for the
  full revision history.)
\end{itemize}

\end{document}
